% AI-Empowered Space-Air-Ground Integrated Networks (SAGIN)
% INFO 3180 Wireless Networks - Group Research Project
% Team: Divyanshu Ratti, Henry Manjo, Rajdeep Kaur

\documentclass[letterpaper]{article} % DO NOT CHANGE THIS
\usepackage[submission]{aaai2026} % DO NOT CHANGE THIS
\usepackage{times} % DO NOT CHANGE THIS
\usepackage{helvet} % DO NOT CHANGE THIS
\usepackage{courier} % DO NOT CHANGE THIS
\usepackage[hyphens]{url} % DO NOT CHANGE THIS
\usepackage{graphicx} % DO NOT CHANGE THIS
\urlstyle{rm} % DO NOT CHANGE THIS
\def\UrlFont{\rm} % DO NOT CHANGE THIS
\usepackage{natbib} % DO NOT CHANGE THIS
\usepackage{caption} % DO NOT CHANGE THIS
\frenchspacing % DO NOT CHANGE THIS
\setlength{\pdfpagewidth}{8.5in} % DO NOT CHANGE THIS
\setlength{\pdfpageheight}{11in} % DO NOT CHANGE THIS

\pdfinfo{
/Title (AI-Empowered Space-Air-Ground Integrated Networks: A Survey)
/Author (Divyanshu Ratti, Henry Manjo, Rajdeep Kaur)
}

\title{AI-Empowered Space-Air-Ground Integrated Networks: A Survey}
\author{
    Divyanshu Ratti \and Henry Manjo \and Rajdeep Kaur \\
    Department of Computer Science / Engineering \\
    INFO 3180 -- Wireless Networks
}

\begin{document}

\maketitle

%%% -------------------------------------------------------
%%% SECTION: Abstract
%%% AUTHOR : Henry Manjo, Divyanshu Ratti
%%% -------------------------------------------------------
\begin{abstract}
Space-Air-Ground Integrated Networks (SAGIN), which combine satellite, aerial, and terrestrial communication layers, are considered a key architecture for future 6G wireless systems because they provide wider coverage and improved connectivity for applications such as disaster recovery, intelligent transportation, and Internet of Things (IoT) services. However, the mobility, heterogeneity, and dynamic nature of SAGIN make these networks difficult to manage using conventional approaches, leading to increasing interest in Artificial Intelligence (AI) techniques such as deep reinforcement learning, federated learning, agentic AI, and edge AI. This paper surveys recent research on AI-enabled SAGIN and reviews how these techniques have been applied to resource allocation, beamforming, trajectory planning, authentication, and network security. Our analysis identifies several common trends across the literature. Most AI-based solutions are validated through simulations rather than real-world deployments, researchers are increasingly adopting decentralized and autonomous network management approaches, and the majority of studies focus on performance optimization while comparatively less attention is given to security. Based on these findings, we discuss current research gaps and identify directions for advancing practical, scalable, and secure AI-enabled SAGIN for future 6G networks.
\end{abstract}

%%% -------------------------------------------------------
%%% SECTION: Introduction
%%% AUTHOR : Rajdeep Kaur
The vision of future 6G wireless communication is to provide fast, reliable, and seamless connectivity everywhere in the world. Unlike traditional wireless networks that mainly depend on terrestrial cell towers, 6G is expected to combine multiple network layers, including ground base stations, drones, high-altitude platforms, and Low Earth Orbit satellites. This type of network is known as a Space-Air-Ground Integrated Network, or SAGIN.

SAGIN is important because it can extend communication coverage to areas where normal ground networks may not be available or reliable. These areas may include remote villages, oceans, mountains, disaster zones, and highly mobile environments. By connecting space, air, and ground communication systems together, SAGIN can support global connectivity and advanced 6G applications such as smart transportation, emergency communication, Internet of Things services, autonomous systems, and military communication.

However, managing SAGIN is challenging because the network is highly dynamic and heterogeneous. Satellites, drones, airborne platforms, and mobile users are constantly moving, which changes signal strength, coverage, and network conditions. These changes can create problems such as signal interference, limited spectrum availability, unstable connections, and difficulty in maintaining smooth communication between different network tiers.

Artificial Intelligence plays an important role in solving these challenges. AI can help manage dynamic spectrum sharing by deciding how radio resources should be assigned between satellites, drones, airborne platforms, and ground networks. AI can also support trajectory optimization by helping drones and airborne platforms choose better movement paths for stronger coverage and lower interference. In addition, AI can improve seamless signal handoffs by predicting when a user should switch from one network layer to another, such as from a terrestrial base station to a drone or satellite connection.

This project focuses on AI-Empowered Space-Air-Ground Integrated Networks for future 6G systems. The main objective is to explore how artificial intelligence can manage spectrum sharing, trajectory optimization, and seamless signal handoffs across space, air, and ground network tiers. The project will also examine related areas such as beamforming, physical layer security, and intelligent network coordination. Through this research, we aim to understand how AI can make SAGIN more efficient, secure, adaptive, and reliable for next-generation wireless communication.

%%% -------------------------------------------------------
%%% SECTION: Background
%%% AUTHOR : Henry Manjo, Divyanshu Ratti, Rajdeep Kaur
%%% -------------------------------------------------------
\section{Background}
% TODO (~0.75 page): Define SAGIN (space/air/ground layers), define
% the relevant AI techniques used across the literature you reviewed
% (e.g., reinforcement learning, federated learning, agentic AI),
% and any key terminology a reader needs before the Literature Review.
Space-Air-Ground Integrated Networks (SAGIN) are communication networks that combine three different layers: space, air, and ground. The space layer includes satellites, the air layer includes platforms such as drones (UAVs) and high-altitude platforms (HAPS), and the ground layer includes cellular base stations, vehicles, IoT devices, and other terrestrial communication infrastructure. By connecting these three layers, SAGIN can provide wider coverage, improve connectivity, and support communication in areas where traditional ground networks may have limited coverage. This integrated architecture is considered an important technology for future 6G networks because it can support applications such as remote sensing, disaster recovery, intelligent transportation, and Internet of Things (IoT) services.

Throughout the literature reviewed in this survey, several recurring AI techniques are used to manage and optimize SAGIN. This section briefly defines each before they are discussed in more detail in the Literature Review.

\textbf{Deep Reinforcement Learning (DRL):} Deep reinforcement learning is an AI technique that combines reinforcement learning with deep neural networks to solve complex decision-making problems. An AI agent learns by interacting with its environment and improving its decisions based on rewards and penalties. In the papers reviewed, DRL was mainly used to optimize UAV trajectory planning, resource allocation, and task offloading while reducing energy consumption and communication delays.

\textbf{Federated Learning (FL):} Federated learning is a machine learning approach where multiple devices train the same AI model without sharing their raw data with a central server. Instead, they only exchange model updates, which helps improve privacy and reduce communication costs. Some of the papers also proposed decentralized federated learning, where satellites communicate directly with each other instead of relying on a central ground station.

\textbf{Agentic AI:} Agentic AI refers to AI systems that can make decisions and perform tasks with minimal human intervention. Instead of only following predefined instructions, these AI agents can monitor network conditions, make decisions, and adapt to changing environments automatically. In future SAGIN and 6G networks, Agentic AI can help improve resource management, routing, and overall network performance.

\textbf{Edge AI:} Edge AI is the use of AI algorithms on devices located close to where data is generated, such as UAVs, satellites, or edge servers, instead of sending all data to a central cloud. This reduces communication delays, lowers bandwidth usage, and allows faster decision-making, which is important for real-time applications in SAGIN.

%%% END SECTION: Background

%%% -------------------------------------------------------
%%% SECTION: Literature Review
%%% AUTHOR : Divyanshu Ratti
%%% -------------------------------------------------------
\section{Literature Review}

Artificial Intelligence (AI) has become one of the main technologies for improving Space-Air-Ground Integrated Networks (SAGIN). Since SAGIN combines satellites, UAVs, high-altitude platforms, and terrestrial networks, it creates a highly dynamic environment where traditional optimization methods often struggle to make decisions quickly. The papers reviewed in this survey show that researchers are increasingly using AI to solve problems related to resource allocation, task offloading, communication optimization, mobility management, security, and intelligent network control. Although different AI techniques are used, most studies have the common goal of making SAGIN more efficient, scalable, and reliable for future 6G communication systems.

\subsection{Resource Allocation and Intelligent Task Offloading}

Resource allocation is one of the most active research areas in AI-enabled SAGIN because communication resources, computing power, and energy are limited across satellites, UAVs, and ground networks. Instead of relying on fixed optimization methods, many researchers use deep reinforcement learning (DRL) to make decisions based on changing network conditions.

Several studies focus on deciding where computing tasks should be processed. Huang et al.~\citep{huang2024joint} proposed a DRL framework that jointly optimizes task offloading and resource allocation across satellites, UAVs, and cloud servers while considering task dependencies. Similarly, Zhou et al.~\citep{Zhou2021IoTTaskScheduling} developed a reinforcement learning scheduler for delay-sensitive IoT applications, allowing UAVs to decide whether tasks should be processed locally or offloaded to nearby edge servers or satellites. Zhang et al.~\citep{zhang2023multiagent} extended this idea by applying multi-agent reinforcement learning to orbital edge computing for the Internet of Remote Things, where multiple AI agents cooperate to improve offloading decisions for remote sensing applications.

Other researchers have focused on improving resource management in specialized environments. Xu and Yu~\citep{xu2025deep} applied deep reinforcement learning to maritime Internet of Things networks, while Zhang et al.~\citep{zhang2025energy} optimized both UAV trajectories and resource allocation for multi-UAV mobile edge computing. Tu et al.~\citep{tu2024priority} also introduced multi-agent reinforcement learning for network slicing, allowing communication resources to be assigned according to different service priorities. Although these studies address different applications, they all show that reinforcement learning is becoming one of the preferred methods for solving complex resource management problems in SAGIN.

\subsection{Communication Optimization, Beamforming, and Mobility}

Besides managing computing resources, many researchers have applied AI to improve wireless communication itself. Beamforming, interference suppression, spectrum sharing, and seamless mobility are important because satellites and UAVs continuously move while wireless conditions change over time.

Khoshkbari et al.~\citep{khoshkbari2025distributed} proposed a distributed beamforming approach where airborne platforms use multi-agent reinforcement learning instead of continuously exchanging channel state information. Mamillapalli et al.~\citep{mamillapalli2026deep} addressed interference suppression by combining deep reinforcement learning with reconfigurable intelligent surfaces (RIS), allowing the network to reduce cross-tier interference while improving spectral efficiency. Kim et al.~\citep{kim2024quantum} further extended cooperative communication by introducing quantum multi-agent reinforcement learning for scheduling CubeSats and high-altitude UAVs.

Mobility management has also received increasing attention. Wu et al.~\citep{wu2025phandover} proposed PHandover, a machine learning-assisted handover framework for LEO satellite networks that predicts signal quality before users switch satellites. Unlike most papers reviewed in this survey, their approach was validated using a real prototype with realistic satellite movement data instead of simulations alone. Similarly, Lin et al.~\citep{lin2025ai} discussed AI-driven seamless and massive network access, showing that intelligent mobility management will become increasingly important as the number of connected users continues to grow.

\subsection{Federated Learning and Decentralized AI}

As SAGIN becomes larger and more distributed, researchers are moving away from centralized AI models toward decentralized learning approaches. Federated learning allows multiple network nodes to train AI models collaboratively without sharing raw data, improving both privacy and communication efficiency.

Yang and Zhang~\citep{yang2024dfedsat} proposed DFedSat, a decentralized federated learning framework where satellites exchange model updates directly through inter-satellite links instead of relying on a central ground station. Qin et al.~\citep{qin2024differentiated} combined federated learning with reinforcement learning for traffic offloading, allowing different regions of the network to maintain local decision-making while still sharing useful learning information. More recently, Li et al.~\citep{li2026adaptive} further improved federated learning by introducing adaptive compressed sensing, differential privacy, and reinforcement learning-based routing to reduce communication overhead while protecting sensitive information.

These studies demonstrate an important trend within recent SAGIN research. Rather than depending on one central controller, future AI systems are expected to distribute intelligence throughout the network, allowing satellites, UAVs, and ground nodes to cooperate more efficiently while reducing communication costs.

\subsection{Generative AI, Semantic Communication, and Intelligent Network Management}

Although reinforcement learning dominates much of the literature, newer AI technologies are beginning to appear in SAGIN research. Generative AI has recently attracted attention because it can generate missing information, improve channel estimation, and support semantic communication.

Zhang et al.~\citep{zhang2024generative} reviewed how generative AI models such as diffusion models, GANs, and transformer-based models can support channel modeling, resource allocation, semantic communication, and network optimization. Huang et al.~\citep{huang2024deep} extended this idea by combining deep reinforcement learning with hybrid bit-level and generative semantic communication, allowing image transmission while reducing bandwidth requirements. Dev et al.~\citep{dev2025advanced} also introduced Agentic AI as a future network management framework where autonomous AI agents continuously monitor network conditions and make intelligent decisions with minimal human intervention. Together, these studies suggest that future SAGIN research will gradually move beyond optimization algorithms toward more autonomous and intelligent network architectures.

\subsection{Security and Future Research Directions}

Compared with resource allocation and communication optimization, relatively fewer studies focus primarily on security. Chou et al.~\citep{chou2024edge} proposed an Edge AI framework for physical layer security, allowing AI models to detect threats closer to where data is generated while reducing communication delay. Yang et al.~\citep{byang2023atmas} introduced an AI-oriented multi-factor authentication framework that continuously verifies users based on spatial and temporal network information rather than relying only on traditional cryptographic authentication.

Several survey papers also provide a broader understanding of the field. Kato et al.~\citep{kato2019optimizing} first demonstrated how AI could improve routing, spectrum management, energy efficiency, and network control in SAGIN, providing one of the earliest foundations for AI-enabled integrated networks. More recent surveys by Bakambekova et al.~\citep{bakambekova2024interplay} and Zhang et al.~\citep{zhang2024management} show that research has expanded rapidly to include reinforcement learning, federated learning, edge AI, generative AI, and autonomous network management. Despite this progress, these surveys also identify common challenges, including limited real-world deployment, scalability, privacy protection, security, and the high computational cost of advanced AI models.

Overall, the literature shows that AI has become an essential technology for enabling future 6G SAGIN. Most current research demonstrates improvements in network performance, resource management, and communication efficiency through simulation-based evaluation. However, relatively few studies validate their approaches in real deployments or fully address security, interoperability, and scalability. These remaining challenges provide important opportunities for future research and motivate the discussion presented in the next section.

%%% END SECTION: Literature Review

%%% -------------------------------------------------------
%%% SECTION: Discussion
%%% AUTHOR : Henry Manjo, Divyanshu Ratti, Rajdeep Kaur
%%% -------------------------------------------------------
\section{Discussion}

The literature reviewed in this survey reveals several recurring patterns and open challenges in applying AI to Space-Air-Ground Integrated Networks (SAGIN).

One pattern noticed across several of the papers is that most of the proposed AI solutions were only evaluated using simulations instead of real-world deployments. Although the simulation results showed improvements in areas such as resource allocation, energy efficiency, communication delay, and network performance, it is still unclear how well these methods would work in real SAGIN environments. Real networks have additional challenges such as satellite mobility, changing communication links, weather conditions, interference, and limited computing resources that are difficult to capture in simulations. This limitation appears in several papers, including the studies on multi-UAV mobile edge computing \citep{zhang2025energy}, decentralized federated learning for LEO satellite constellations \citep{yang2024dfedsat}, and deep reinforcement learning for task offloading and resource allocation \citep{huang2024joint}. This suggests that while AI-based approaches show a lot of promise, more real-world testing is needed before they can be widely deployed in future SAGIN systems.

Another pattern noticed is that many of the papers are moving away from centralized network management and instead using more decentralized or autonomous approaches. For example, the DFedSat paper removes the need for a central ground station by allowing satellites to train AI models directly with each other \citep{yang2024dfedsat}. The distributed beamforming paper uses multi-agent reinforcement learning so that each airborne platform can make decisions without constantly sharing channel state information \citep{khoshkbari2025distributed}. The Agentic AI paper also follows this trend by proposing intelligent AI agents that can monitor network conditions and make decisions automatically instead of relying on human operators \citep{dev2025advanced}. Even though these approaches use different AI techniques, they all aim to reduce communication overhead, improve scalability, and make the network more adaptable. However, they also introduce new challenges, such as coordinating many independent agents and ensuring that decentralized systems remain reliable as the network grows.

A third pattern noticed is that most of the papers focus on improving network performance rather than security. Many of the studies use AI techniques such as deep reinforcement learning, multi-agent reinforcement learning, or Agentic AI to optimize resource allocation, beamforming, task offloading, and communication efficiency. In comparison, only a few papers, such as Chou's work on Edge AI-enabled physical layer security \citep{chou2024edge}, focus mainly on protecting the network from security threats. The survey paper by Bakambekova et al. also points out that security remains one of the major open challenges in AI-enabled SAGIN \citep{bakambekova2024interplay}. This suggests that although researchers are making good progress in improving network performance, more work is still needed to develop AI solutions that can provide both high performance and strong security at the same time.

%%% END SECTION: Discussion

%%% -------------------------------------------------------
%%% SECTION: Proposed Idea
%%% AUTHOR : Divyanshu Ratti
%%% -------------------------------------------------------
\section{Proposed Idea}

Based on the literature reviewed in this survey, we believe that future research should focus on developing a more integrated AI framework instead of using individual AI techniques to solve separate networking problems. Most of the existing studies apply deep reinforcement learning, federated learning, or multi-agent reinforcement learning to improve one specific task, such as resource allocation, task offloading, or beamforming. Although these approaches show promising results, they are usually evaluated only through simulations and often operate independently from one another \citep{huang2024joint,qin2024differentiated,zhang2024management}.

To address these limitations, we propose a hierarchical AI framework where satellites, UAVs, and ground stations each contain a local AI agent that can monitor network conditions and make decisions independently. These agents can collaborate using federated learning so that they improve a shared AI model without exchanging sensitive user data \citep{yang2024dfedsat,li2026adaptive}. Multi-agent reinforcement learning can then be used to coordinate decisions between different network layers, allowing the system to respond more efficiently to changing traffic, mobility, and communication conditions \citep{kim2024quantum,zhang2023multiagent}.

To further improve network management, Generative AI could be added as a decision support tool for network operators. Instead of directly controlling the network, it could analyse historical network data, predict possible congestion, and recommend optimization strategies before performance is affected \citep{zhang2024generative}. This could help network administrators make faster and more informed decisions while reducing manual network management.

Although this framework would require further research and real-world validation, it combines several AI techniques that have already shown promising results individually. We believe that integrating these methods into a single framework could improve scalability, adaptability, and privacy while supporting the development of more intelligent and reliable SAGIN for future 6G wireless networks.

%%% END SECTION: Proposed Idea

%%% -------------------------------------------------------
%%% SECTION: Conclusion
%%% AUTHOR : Divyanshu Ratti, Rajdeep Kaur
%%% -------------------------------------------------------
\section{Conclusion}
% TODO (~0.25 page): Summarize the survey and restate the proposed
% direction / future work in 3-4 sentences.
Placeholder conclusion text.

%%% END SECTION: Conclusion

%%% -------------------------------------------------------
%%% SECTION: References
%%% AUTHOR : Henry Manjo, Divyanshu Ratti, Rajdeep Kaur
%%% -------------------------------------------------------
% References are pulled from references.bib using natbib.
% Add entries to references.bib, then cite them above with \citep{key}.
\bibliography{references}
\bibliographystyle{aaai2026}

%%% END SECTION: References

\end{document}
